\section*{\centering{RESUMO}} 
\noindent A Dívida Técnica (DT) de documentação é um desafio na indústria de \textit{software}. No Sistema de Fomento ao Microempreendedor (SFM), o fluxo de desenvolvimento empírico resultou no acúmulo de DT, gerando falhas de comunicação e inviabilizando testes. Este trabalho objetiva relatar a análise dessa DT e propor sua mitigação mediante artefatos formais. A metodologia englobou elicitação de requisitos, construção de uma Matriz de Rastreabilidade e aplicação de questionário a sete desenvolvedores para aferir a percepção de completude do sistema. Os resultados revelaram que 82\% dos fluxos não apresentavam cobertura de teste. Ademais, o questionário refutou a observação empírica inicial da pesquisa: um módulo previamente avaliado como consolidado revelou, na visão da equipe, implementação parcial e alta divergência de compreensão. Constatou-se, ainda, que os colaboradores reconhecem o alto potencial da documentação produzida para auxiliá-los em suas atividades diárias. Conclui-se que a documentação formal fornece uma base estrutural sólida, capaz de mitigar a opacidade diagnosticada, orientar a codificação e viabilizar a futura automação de testes, atestando a importância da Engenharia de Requisitos para a qualidade do \textit{software}

\noindent {\bf Palavras-chave:} $<$Dívida Técnica$>$, $<$Engenharia de Requisitos$>$, $<$Engenharia de Software$>$, $<$Matriz de Rastreabilidade$>$, $<$Teste de Software$>$.
